% LSPV-4 REPORT::READCEIL1::2026-10-12T17:57:24Z
% !TEX program = xelatex
\documentclass[11pt]{article}
\usepackage{fontspec}
\usepackage[margin=0.86in]{geometry}
\usepackage{amsmath,amssymb,amsthm,booktabs,longtable,tabularx,array,enumitem,xurl,hyperref,microtype}
\hypersetup{colorlinks=true,linkcolor=blue,urlcolor=blue,citecolor=blue}
\newtheorem{theorem}{Theorem}
\newtheorem{corollary}{Corollary}
\newtheorem{definition}{Definition}
\newcommand{\SAT}{\mathsf{SAT}}
\newcommand{\Core}{\mathcal C_{2,2}^{\pm}}
\newcommand{\EXACT}{\textbf{PROVED/EXACT}}
\newcommand{\KNOWN}{\textbf{KNOWN THEOREM}}
\newcommand{\COMP}{\textbf{COMPUTATIONAL}}
\newcommand{\COND}{\textbf{CONDITIONAL}}
\newcommand{\CAND}{\textbf{CANDIDATE}}
\newcommand{\OPEN}{\textbf{OPEN}}
\newcommand{\RET}{\textbf{FALSIFIED/RETIRED}}
\newcommand{\same}{\textsc{same\_realization}}

\title{READCEIL1: Actual reads universally certify the decision,\\
but their size has a linear ceiling}
\author{LSPV-4 / UNIVERSAL\_CORE\_BRIDGE}
\date{12 October 2026, 17:57 UTC}

\begin{document}
\maketitle

\begin{center}
\fbox{\parbox{0.94\linewidth}{
\textbf{Executive status.} P versus NP remains \OPEN.  For every deterministic
decider and input, the distinct input positions and values actually read form
a monochromatic decision certificate.  This is \EXACT{} and is the standard
decision-tree certificate principle.  Its cardinality is nevertheless at most
the encoded input length \(N\), and the whole input can be transported in
\(O(N)\) time.  Replaying the run retains the decider's own polynomial
exponent.  Thus \texttt{READ-CERTIFICATE-SIZE-AS-SUPERPOLY-BRIDGE1} is
\RET.  The direct target remains a superpolynomial lower bound for every
uniform deterministic decider of the fixed NP-complete clean core; no
extension to unrestricted SAT formula families is required.}}
\end{center}

\section{Claim card and exact target}
\begin{description}[leftmargin=0cm,labelwidth=0cm]
\item[Claim ID.] \texttt{READ-CERTIFICATE-CEILING1}.
\item[Status before.] \CAND.  The trace program asked whether a data-bearing
object weaker than full residual compatibility is forced by every decision.
\item[Desired conclusion.] Find a universal same-execution object and test
whether its size, extraction, reconstruction, or replay cost can force a
superpolynomial lower bound.
\item[Input family.] Canonically encoded signed Monotone (3)SAT-((2,2)),
denoted \(\Core\), with \(N=|\langle F\rangle|\).  Invalid encodings are
rejected by the ordinary total language decider.
\item[Algorithm quantifier.] Every total uniform deterministic machine \(A\)
deciding \(\Core\); each machine may have its own polynomial exponent.
\item[Model.] A deterministic random-access input machine with all address
computation, preprocessing, work, output and space charged.  A Turing input
tape is covered by treating visited cells as read positions.
\item[Resource.] Distinct input position/value pairs read before halt, not a
full timed transcript.
\item[\same{} objects.] The read pairs arise in the one run \(A(x)\).  Any
later replay or compatibility computation is separately charged unless it is
explicitly embedded.
\item[Direct endpoint.] It would suffice to prove
\[
 \forall A\in\operatorname{Dec}(\Core)\ \forall d\ \forall N_0\
 \exists F\in\Core:\ |F|\ge N_0\ \wedge\ T_A(F)>|F|^d.       \tag{1}
\]
Since \(\Core\) is NP-complete, (1) implies \(\mathsf P\ne\mathsf{NP}\)
without broadening the formula family.  Equation (1) is \OPEN, not assumed.
\item[Falsifier.] A universal upper bound \(|Q_A(x)|\le N\), or acquisition by
plain input transport, defeats any superpolynomial conclusion based only on
certificate size.
\end{description}

\section{Predecessor recovery and repository state}
The cyclic predecessor is LSPV-3.  Its newest canonical report,
\path{research/agents/lspv-3/2026-10-13_0236_LSPV-3_REVUNCOMP1.tex}, was read
from GitHub and is substantive, with the required header and sentinel.  No
recovery request or duplicate file is needed.

The standing repair request for MODELBRIDGE-0022 is resolved:
\path{research/agents/lspv-4/2026-10-12_1002_LSPV-4_MODELBRIDGE-0022.tex}
exists canonically and verifies.  No missing scientific body was reconstructed.

The two current obstacles are (i) \emph{universality}---the object must exist
for every decider, and (ii) \emph{strength}---it must exceed linear source
transport or otherwise lower-bound the computation producing the answer.  The
first is solved below; the second fails for certificate cardinality.

\section{Read/value certificate theorem}
\begin{definition}[Actual-read certificate]
Fix a deterministic machine \(A\), an input length \(N\), and
\(x\in\{0,1\}^N\).  Let
\[
 R_A(x)=\{i\in[N]:A(x)\text{ reads position }i\text{ before halting}\},
\quad
 Q_A(x)=\{(i,x_i):i\in R_A(x)\}.                            \tag{2}
\]
Repeated reads appear once in \(Q_A(x)\).
\end{definition}

\begin{theorem}[\texttt{READ-CERTIFICATE-CEILING1}]
For every total deterministic machine \(A\), every \(N\), and every
\(x\in\{0,1\}^N\):
\begin{enumerate}[label=(\roman*),leftmargin=1.7em]
\item if \(y_i=x_i\) for every \(i\in R_A(x)\), then \(A(x)\) and \(A(y)\)
have the same sequence of internal states, requested addresses, read values,
halting time and output;
\item \(Q_A(x)\) is a monochromatic certificate for the Boolean function
computed by \(A\) on length \(N\);
\item \(|Q_A(x)|=|R_A(x)|\le N\).
\end{enumerate}
\end{theorem}

\begin{proof}
Induct on transition number.  Initially the internal configurations coincide.
Equal configurations yield the same next operation and, on an input read, the
same address \(i\).  That address belongs to \(R_A(x)\), so the hypothesis
gives \(y_i=x_i\); the next configurations remain equal.  Continue through
halting to prove (i).  Equal output on the cylinder defined by (2) proves (ii).
Since \(R_A(x)\subseteq[N]\), (iii) follows.
\end{proof}

This is unconditional and exact.  It matches the standard certificate notion
in deterministic decision-tree complexity~\cite{BdW02}; the proof records the
adaptive-address machine form needed here.

\begin{corollary}[Linear ceiling]
For every language, including \(\Core\), a lower bound charging only the
number of distinct read/value pairs is at most \(N\).  The full certificate
\[
 Q_{\rm full}(x)=\{(i,x_i):i\in[N]\}                        \tag{3}
\]
is acquired by a fixed scan in \(O(N)\) time and \(O(N)\) output space, or
\(O(\log N)\) writable workspace when streamed.
\end{corollary}

Thus actual reads are a genuine universal decision object but not a
superpolynomial bottleneck.  Exponentially many semantic residual questions
may be definable from the source while (3) still contains only \(N\) bits.

\section{Replay and exponent accounting}
Assume counterfactually that \(\Core\in\mathsf P\) and \(A\) runs in
\(T_A(N)\le N^{K_A}\), \(S_A(N)\le N^{P_A}\).  Recording distinct reads costs
at most
\[
 T_{\rm record}(N)=T_A(N)+N^{e+o(1)},\qquad
 S_{\rm record}(N)=S_A(N)+N^{p+o(1)}.                       \tag{4}
\]
Replaying \(A\) to validate the certificate costs
\[
 T_{\rm replay}(N)=N^{\max\{K_A,r,e\}+o(1)}.                \tag{5}
\]
These are upper-bound accounting identities, not lower bounds.

For each fixed proposed universal exponent \(c\), the extensionally equal
wrapper \(A_d\) may perform \(N^d\) input-oblivious steps before running \(A\).
A literal same-run record or replay has exponent at least \(d\).  Optimizing
away the delay changes the measured realization and returns to an external
source computation.  This exponent/\same{} guard is adopted from earlier
black-box and normalization reports rather than claimed as new.

Consequently, charging only cardinality gives the ceiling \(N\); charging
actual generation/replay retains \(K_A\); and quotienting implementation cost
turns the object into source transport.  This retires
\texttt{READ-CERTIFICATE-SIZE-AS-SUPERPOLY-BRIDGE1}.

\section{Bounded reproducible audit}
The script \path{experiments/review/read_certificate_ceiling_bridge.py}
writes \path{results/raw/review/2026-10-12_read_certificate_ceiling.json}.
It completed in (3.035) seconds with (298{,}132) assertions and zero
failures; no chunking was needed.

The audit exhausts all Boolean functions for \(n\le3\), checks 4,096 seeded
four-bit functions, and evaluates every input.  Its canonical adaptive tree
queries coordinates in order until the current cylinder is monochromatic.
Every returned read set certified the output on every compatible completion;
no certificate exceeded \(n\), and maxima were \(1,2,3,4\).

Twelve actual CNFs were brute-force checked: satisfiable and unsatisfiable
2SAT, Horn, affine/Tseitin, bounded-treewidth, independent-product,
padded/redundant, and alternate-encoding controls.  DIRECT, COPY, DEAD, DELAY,
INTERP, and RECOMP gave 72 answer-preserving cells.  A \(24^3\) exponent grid
checked 41,472 inequalities for \(\max\{K_A,e,r\}\).  The canonical payload
SHA-256 is
\path{33f361d944a3a58b755ea311a3c6f3dfd5e7806804a393ab70e74e5c3857fa3d}.
These tests validate identities and bookkeeping only.

\section{Compatibility and the repaired target}
For \(F\in\Core\) and partial assignment \(\alpha\), retain
\[
 H(F,\alpha)=\{\beta:\alpha\cup\beta\models F\},\qquad
 h(F,\alpha)=\mathbf 1[H(F,\alpha)\ne\varnothing].          \tag{6}
\]
The theorem says that the cylinder consistent with \(Q_A(F)\) has one decision
label.  It does not say \(Q_A(F)\) explicitly represents or efficiently
evaluates all values (6).  Even \(Q_{\rm full}(F)\) merely reproduces the
source.  Enumerating all of \(H\) has exponential output, while evaluating one
\(h(F,\alpha)\) is a separate SAT decision after restriction and may require a
charged reduction back into the fixed core.

An earlier READCERT1 concept was reported only as a local package outside the
canonical repository; its exact source body was unavailable.  It was not
reconstructed or treated as verified.  The present theorem was derived
independently from the canonical corpus and standard certificate principle.
It deliberately avoids any promise-domain assertion that a total language
decider can skip syntax validation on malformed encodings.

The repaired target is
\begin{quote}
\texttt{DECISION-SELECTION-WORK1} (\OPEN): for every uniform deterministic
decider \(A\) of the fixed NP-complete core and every exponent \(d\), find
arbitrarily long core inputs on which the fully charged computation selecting
the correct one-bit label exceeds \(N^d\), while allowing full input reads,
preprocessing, recomputation, affine elimination, compact representations and
equivalent wrappers.
\end{quote}
This is the direct non-\(P\) target, not an asserted intermediate theorem.

\section{Faithful main4/main5 comparison and peer reconciliation}
The actual \path{main4.tex} is the ARC achievable-region manuscript and
\path{main5.tex} is the WRAPPER construction manuscript.  Both preserve
all-input acquisition/coverage and P versus NP as OPEN.  Neither asserts that
every clean-core decider uses its DAG, compatibility representation, proof
system, or wrapper.  This report supplies a weaker universal interface and
proves why its cardinality is insufficient.

INFONOGO1 proves that mutual information with deterministic \(H=g(X)\)
measures source correlation, not provenance; COPY saturates it cheaply.
INFOFLOWCHAIN1 shows serial transport realizes per-step information increments.
REVUNCOMP1 shows reversible compute--copy--uncompute removes generic logical
erasure.  STATE\_READBACK1 proves an exact service-level rank bound but not a
service obligation for ordinary deciders.  TRACECONJ1 combines source
injectivity, fixed schedules, causal liveness and positive work without forcing
\(H\)-specific semantics.  READCEIL1 differs positively because the read set is
unconditionally forced; its linear ceiling then blocks the desired conclusion.

\section{Literature and barrier-transfer audit}
Buhrman--de Wolf provide the standard query/certificate framework~\cite{BdW02}:
\textbf{DIRECTLY TRANSFERABLE} at that level, not as unrestricted time.
Cook's NP-completeness framework makes the scope correction exact: placing the
fixed NP-complete core outside \(P\) suffices~\cite{CookPvsNP}.  Döcker and
Darmann--Döcker establish the relevant bounded-occurrence monotone family
~\cite{Doecker19,DarmannDoecker21}.  These are \KNOWN{} results, not lower
bounds.

Baker--Gill--Solovay and Aaronson--Wigderson remain barriers to relativizing or
algebrizing approaches~\cite{BGS75,AW09}.  The read-certificate theorem
relativizes and uses no nonalgebrizing step.  Razborov--Rudich is not directly
triggered because no large constructive circuit property is proved
~\cite{RR97}.  FLMV applies only to its stated time--space region
~\cite{FLMV05}.  VLSI, communication, streaming, proof-system, and bounded-read
service lower bounds are \textbf{FAILED TRANSFERS} without a matched machine
and interface.  Blum's axiomatic measures support explicit program-cost
accounting but do not force a chosen resource: \textbf{USEFUL ANALOGY needing
a new transfer theorem}~\cite{Blum67}.

\section{Post-quantum technique ledger}
\begin{longtable}{p{0.28\linewidth}p{0.65\linewidth}}
\toprule Technique & Classification and consequence\\ \midrule\endhead
Schur powers, puncturing, shortening & \textbf{FAILED TRANSFER}: a read set is
not a forced linear code whose Schur square controls decision time.\\
Quadratic relations, Hilbert functions, syzygies & \textbf{FAILED TRANSFER}
without a forced algebraic representation.\\
Lossiness/collision entropy & \textbf{EXACT OBJECT/THEOREM TRANSFER} for the
map from inputs to read certificates: fibers are compatible cylinders.  Large
fibers certify a common answer and can help early stopping; they do not imply
hardness.\\
Worst-case/average-case architecture & \textbf{ANALOGY/CANDIDATE}: external
extraction has the right direction, but all calls and solver exponents remain
charged.\\
Collision extraction / special soundness & \textbf{FAILED TRANSFER}: two
inputs sharing a certificate are different source statements, not conflicting
openings of one commitment.\\
Hash/random oracle & \textbf{FAILED TRANSFER} unless every assumption is
removed.\\ \bottomrule
\end{longtable}

The structure is algorithmically exploitable: a small certificate permits
early stopping.  It is not a hardness source.

\section{Relativization, nonuniformity, and one-bit output}
The induction survives oracle access and low-degree extensions; it supplies no
barrier escape.  Nonuniform advice cannot enlarge the read-set ceiling beyond
\(N\), while the target remains uniform.  Compact circuits, branching programs,
BDDs and d-DNNFs may summarize semantics but their lower bounds do not transfer
automatically to RAM/Turing computation.  A decision machine owes one bit.  It
does not owe a witness, all residual answers, a compatibility matrix, or a
canonical proof without a solver-independent reduction.

\section{CORPUS REVIEWED}
This run surveyed the GitHub tree and read root \path{main.tex},
\path{AGENTS.md}, \path{CODEX.md}, the complete PNP skill/protocol and linked
mode files, all \path{tasks/}, actual \path{main4.tex}/\path{main5.tex}, newest
verified MAIN TRACECONJ1, and latest canonical LSPV-1--4 reports INFONOGO1,
INFOFLOWCHAIN1, REVUNCOMP1, and STATE\_READBACK1.  It inspected historical
black-box extraction, witness complexity, implementation replacement,
two-axis covariance, normalized cuts, full-trace source leakage, and
schedule-plus-decision reports; the HEC/VSC-HEC/G-PPC-HEC/CRS-PPC/CCSPC/RCCN
branch; governance; and cited experiments.  Search covered read certificates,
decision trees, trace projections, self-reduction, normalization, oracle
extensionality, padding, and wrappers.  Primary literature was searched before
promotion.

\section{DUPLICATION / CONTRADICTION AUDIT}
The certificate principle is standard and appeared in older project material;
it is not claimed as a new mathematical discovery.  The contribution is its
exact integration with the current core target: adaptive-read proof, linear
ceiling, full-input control, exponent ledger, and retirement boundary.
READBACK1 concerns a declared post-hoc service; (2) records source positions
the original run read.  The information-flow peers predict that source
transport can be cheap.  No main4/main5 construction claim is contradicted.

The local-only READCERT1 description included a promise-domain short-circuit
example whose behavior on invalid encodings was unavailable.  This report does
not adopt it and preserves total-language syntax validation.

\section{WHAT CHANGED RELATIVE TO ROOT MAIN}
Root MAIN already states that source reads and source-injective traces do not
force compatibility computation.  This run adds the positive universal
certificate theorem, closes the cardinality route at \(N\), records exponent
survival for replay, and repairs the target to decision-selection work on the
fixed core rather than transfer to broader SAT inputs.

\section{ADOPTED / REJECTED / RETIRED PRIOR COMPONENTS}
\textbf{Adopted:} fixed NP-complete signed core; direct-core scope;
deterministic certificates; one-bit output; full accounting; \same{}; hostile
wrappers; and STATE\_READBACK1 within its service model.

\textbf{Rejected as universal transfers:} certificate size to unrestricted
time; full source access to computation of \(H\); replay cost to original-run
cost; proof-system, circuit, branching-program, streaming, VLSI, or
communication lower bounds to arbitrary clean-core deciders.

\textbf{FALSIFIED/RETIRED:}
\texttt{READ-CERTIFICATE-SIZE-AS-SUPERPOLY-BRIDGE1}, alongside already retired
bare full-trace factorization, schedule-plus-decision compatibility recovery,
mutual-information provenance, per-step information cost, generic erasure cost,
and unqualified readback-service transfer.

\section{STALE OR OVERSTATED CLAIMS FOUND}
The claim that a proved lower bound for the fixed NP-complete core must still be
extended to unrestricted SAT inputs is stale.  Universality is over algorithms
deciding the core.  A read certificate cannot be superlinear in distinct input
positions.  Reconstructing the source from a trace does not show that the
solver computed residual compatibility; it starts a separate charged task.

\section{Strict status and handoff}
\begin{center}
\begin{tabularx}{0.98\linewidth}{Xp{0.22\linewidth}X}
\toprule Item & Status & Reason\\ \midrule
Actual-read certificate theorem & \EXACT & Complete induction proof\\
Linear certificate ceiling & \EXACT & Read positions lie in ([N])\\
Finite Boolean/CNF/exponent audit & \COMP & 298,132 assertions, zero failures\\
Read-certificate size as superpolynomial bridge & \RET & Universal bound \(N\)\\
Decision-selection work lower bound on fixed core & \OPEN & Direct target\\
Model-matched unconditional contradiction & \OPEN & General-machine LB missing\\
P versus NP & \OPEN & Load-bearing arrow open\\ \bottomrule
\end{tabularx}
\end{center}

LSPV-1 should seek an exponent-free obligation to select the decision after
the full source is read.  LSPV-2 should attack it with streaming, Gaussian
elimination, bounded-treewidth DP, compact circuits and products.  LSPV-3
should apply DEAD/COPY/DELAY/INTERP/RECOMP and reversible uncomputation.
LSPV-5 should remove broader-input transfer from the direct route while
retaining all-algorithm and model-matching requirements.

\textbf{Highest-value next target.} Prove (1), or identify an exponent-free
one-bit-decision obligation stronger than source observation but weaker than
full compatibility, and prove every core decider incurs it in the same
realization.  It must survive complete input reads and must not restate (1).

\begin{thebibliography}{99}
\bibitem{BdW02} H. Buhrman and R. de Wolf, ``Complexity measures and decision
tree complexity: a survey,'' \emph{TCS} 288(1), 21--43, 2002,
\url{https://homepages.cwi.nl/~rdewolf/publ/qc/dectree.pdf}.
\bibitem{CookPvsNP} S. Cook, ``The P versus NP Problem,'' Clay Mathematics
Institute, \url{https://www.claymath.org/wp-content/uploads/2022/06/pvsnp.pdf}.
\bibitem{Doecker19} J. Döcker, ``Monotone 3-Sat-(2,2) is NP-complete,''
arXiv:1912.08032, \url{https://arxiv.org/abs/1912.08032}.
\bibitem{DarmannDoecker21} A. Darmann and J. Döcker, ``On simplified
NP-complete variants of Monotone 3-Sat,'' \emph{Discrete Applied Mathematics}
292, 45--58, 2021, \url{https://doi.org/10.1016/j.dam.2020.12.010}.
\bibitem{BGS75} T. Baker, J. Gill, and R. Solovay, ``Relativizations of the
P=?NP question,'' \emph{SIAM J. Comput.} 4(4), 431--442, 1975,
\url{https://doi.org/10.1137/0204037}.
\bibitem{AW09} S. Aaronson and A. Wigderson, ``Algebrization: a new barrier in
complexity theory,'' \emph{ACM TOCT} 1(1), 2009,
\url{https://doi.org/10.1145/1490270.1490272}.
\bibitem{RR97} A. Razborov and S. Rudich, ``Natural proofs,'' \emph{JCSS}
55(1), 24--35, 1997, \url{https://doi.org/10.1006/jcss.1997.1494}.
\bibitem{FLMV05} L. Fortnow, R. Lipton, D. van Melkebeek, and A. Viglas,
``Time-space lower bounds for satisfiability,'' \emph{JACM} 52(6), 835--865,
2005, \url{https://doi.org/10.1145/1101821.1101822}.
\bibitem{Blum67} M. Blum, ``A machine-independent theory of the complexity of
recursive functions,'' \emph{JACM} 14(2), 322--336, 1967,
\url{https://doi.org/10.1145/321386.321395}.
\end{thebibliography}

% END_OF_VERIFIED_REPORT::LSPV-4::READCEIL1
\end{document}
